The term geostationary (GEO) refers to a satellite whose orbit is on the equatorial plane, its orbit is
circular, and its orbital period is equal to Earth's period of revolution around its axis.
For this problem, all orbits around Earth are considered circular. An important characteristic of the
orbit of a satellite is the elevation angle of the satellite, which is defined as the angle at which the
satellite is visible above the local horizon for an observer on Earth.
\begin{description}
\item[(A)] What is the maximum latitude the geostationary satellite is visible at? \hfill\mbox{[12 points]}
\end{description}
Geostationary satellites, because of their distance from the center of the Earth, observe only a fraction
of the Earth's surface. If a satellite needs to view a larger fraction of the Earth, it must change its orbit using rockets.
\begin{description}
\item[(B)] The satellite moves from its location and observes the Earth having an angular diameter of $10^\circ$.
What is the distance of the satellite from the center of the Earth? \hfill\mbox{[10 points]}
\item[(C)] How much does the total (kinetic and potential) energy of the satellite need to change to move
from the geostationary orbit to the one described in question (B)? The satellite has a mass of 1000 kg. \hfill\mbox{[10 points]}
\item[(D)] What is the ratio of the gravitational acceleration $g$ on the Earth's surface compared to the
gravitational acceleration at the location described in question (B)? \hfill\mbox{[8 points]}
\end{description}
\textit{Data from the accompanying constants sheet.}
$G=6.674\times10^{-11}\,\mathrm{N\,m^2\,kg^{-2}}$, $M_E=5.972\times10^{24}\,\mathrm{kg}$,
$R_E=6378$ km, and one sidereal day is $23^{\mathrm h}56^{\mathrm m}04^{\mathrm s}$.
